% Options for packages loaded elsewhere
\PassOptionsToPackage{unicode}{hyperref}
\PassOptionsToPackage{hyphens}{url}
\PassOptionsToPackage{dvipsnames,svgnames,x11names}{xcolor}
%
\documentclass[
]{article}
\usepackage{amsmath,amssymb}
\usepackage{iftex}
\ifPDFTeX
  \usepackage[T1]{fontenc}
  \usepackage[utf8]{inputenc}
  \usepackage{textcomp} % provide euro and other symbols
\else % if luatex or xetex
  \usepackage{unicode-math} % this also loads fontspec
  \defaultfontfeatures{Scale=MatchLowercase}
  \defaultfontfeatures[\rmfamily]{Ligatures=TeX,Scale=1}
\fi
\usepackage{lmodern}
\ifPDFTeX\else
  % xetex/luatex font selection
  \setmonofont[Scale=0.85]{DejaVu Sans Mono}
\fi
% Use upquote if available, for straight quotes in verbatim environments
\IfFileExists{upquote.sty}{\usepackage{upquote}}{}
\IfFileExists{microtype.sty}{% use microtype if available
  \usepackage[]{microtype}
  \UseMicrotypeSet[protrusion]{basicmath} % disable protrusion for tt fonts
}{}
\makeatletter
\@ifundefined{KOMAClassName}{% if non-KOMA class
  \IfFileExists{parskip.sty}{%
    \usepackage{parskip}
  }{% else
    \setlength{\parindent}{0pt}
    \setlength{\parskip}{6pt plus 2pt minus 1pt}}
}{% if KOMA class
  \KOMAoptions{parskip=half}}
\makeatother
\usepackage{xcolor}
\usepackage[a4paper,margin=2.2cm]{geometry}
\usepackage{color}
\usepackage{fancyvrb}
\newcommand{\VerbBar}{|}
\newcommand{\VERB}{\Verb[commandchars=\\\{\}]}
\DefineVerbatimEnvironment{Highlighting}{Verbatim}{commandchars=\\\{\}}
% Add ',fontsize=\small' for more characters per line
\newenvironment{Shaded}{}{}
\newcommand{\AlertTok}[1]{\textcolor[rgb]{1.00,0.00,0.00}{\textbf{#1}}}
\newcommand{\AnnotationTok}[1]{\textcolor[rgb]{0.38,0.63,0.69}{\textbf{\textit{#1}}}}
\newcommand{\AttributeTok}[1]{\textcolor[rgb]{0.49,0.56,0.16}{#1}}
\newcommand{\BaseNTok}[1]{\textcolor[rgb]{0.25,0.63,0.44}{#1}}
\newcommand{\BuiltInTok}[1]{\textcolor[rgb]{0.00,0.50,0.00}{#1}}
\newcommand{\CharTok}[1]{\textcolor[rgb]{0.25,0.44,0.63}{#1}}
\newcommand{\CommentTok}[1]{\textcolor[rgb]{0.38,0.63,0.69}{\textit{#1}}}
\newcommand{\CommentVarTok}[1]{\textcolor[rgb]{0.38,0.63,0.69}{\textbf{\textit{#1}}}}
\newcommand{\ConstantTok}[1]{\textcolor[rgb]{0.53,0.00,0.00}{#1}}
\newcommand{\ControlFlowTok}[1]{\textcolor[rgb]{0.00,0.44,0.13}{\textbf{#1}}}
\newcommand{\DataTypeTok}[1]{\textcolor[rgb]{0.56,0.13,0.00}{#1}}
\newcommand{\DecValTok}[1]{\textcolor[rgb]{0.25,0.63,0.44}{#1}}
\newcommand{\DocumentationTok}[1]{\textcolor[rgb]{0.73,0.13,0.13}{\textit{#1}}}
\newcommand{\ErrorTok}[1]{\textcolor[rgb]{1.00,0.00,0.00}{\textbf{#1}}}
\newcommand{\ExtensionTok}[1]{#1}
\newcommand{\FloatTok}[1]{\textcolor[rgb]{0.25,0.63,0.44}{#1}}
\newcommand{\FunctionTok}[1]{\textcolor[rgb]{0.02,0.16,0.49}{#1}}
\newcommand{\ImportTok}[1]{\textcolor[rgb]{0.00,0.50,0.00}{\textbf{#1}}}
\newcommand{\InformationTok}[1]{\textcolor[rgb]{0.38,0.63,0.69}{\textbf{\textit{#1}}}}
\newcommand{\KeywordTok}[1]{\textcolor[rgb]{0.00,0.44,0.13}{\textbf{#1}}}
\newcommand{\NormalTok}[1]{#1}
\newcommand{\OperatorTok}[1]{\textcolor[rgb]{0.40,0.40,0.40}{#1}}
\newcommand{\OtherTok}[1]{\textcolor[rgb]{0.00,0.44,0.13}{#1}}
\newcommand{\PreprocessorTok}[1]{\textcolor[rgb]{0.74,0.48,0.00}{#1}}
\newcommand{\RegionMarkerTok}[1]{#1}
\newcommand{\SpecialCharTok}[1]{\textcolor[rgb]{0.25,0.44,0.63}{#1}}
\newcommand{\SpecialStringTok}[1]{\textcolor[rgb]{0.73,0.40,0.53}{#1}}
\newcommand{\StringTok}[1]{\textcolor[rgb]{0.25,0.44,0.63}{#1}}
\newcommand{\VariableTok}[1]{\textcolor[rgb]{0.10,0.09,0.49}{#1}}
\newcommand{\VerbatimStringTok}[1]{\textcolor[rgb]{0.25,0.44,0.63}{#1}}
\newcommand{\WarningTok}[1]{\textcolor[rgb]{0.38,0.63,0.69}{\textbf{\textit{#1}}}}
\usepackage{longtable,booktabs,array}
\usepackage{calc} % for calculating minipage widths
% Correct order of tables after \paragraph or \subparagraph
\usepackage{etoolbox}
\makeatletter
\patchcmd\longtable{\par}{\if@noskipsec\mbox{}\fi\par}{}{}
\makeatother
% Allow footnotes in longtable head/foot
\IfFileExists{footnotehyper.sty}{\usepackage{footnotehyper}}{\usepackage{footnote}}
\makesavenoteenv{longtable}
\setlength{\emergencystretch}{3em} % prevent overfull lines
\providecommand{\tightlist}{%
  \setlength{\itemsep}{0pt}\setlength{\parskip}{0pt}}
\setcounter{secnumdepth}{5}
\ifLuaTeX
\usepackage[bidi=basic]{babel}
\else
\usepackage[bidi=default]{babel}
\fi
\babelprovide[main,import]{english}
% get rid of language-specific shorthands (see #6817):
\let\LanguageShortHands\languageshorthands
\def\languageshorthands#1{}
\ifLuaTeX
  \usepackage{selnolig}  % disable illegal ligatures
\fi
\usepackage{bookmark}
\IfFileExists{xurl.sty}{\usepackage{xurl}}{} % add URL line breaks if available
\urlstyle{same}
\hypersetup{
  pdftitle={VERA X16 for Atari --- Running the tests in the emulator},
  pdflang={en},
  colorlinks=true,
  linkcolor={blue},
  filecolor={Maroon},
  citecolor={Blue},
  urlcolor={Blue},
  pdfcreator={LaTeX via pandoc}}

\title{VERA X16 for Atari --- Running the tests in the emulator}
\usepackage{etoolbox}
\makeatletter
\providecommand{\subtitle}[1]{% add subtitle to \maketitle
  \apptocmd{\@title}{\par {\large #1 \par}}{}{}
}
\makeatother
\subtitle{atari800 fork with the VeraX16 PBI card}
\author{Gianluca Renzi}
\date{2026-10-06}

\begin{document}
\maketitle

{
\hypersetup{linkcolor=}
\setcounter{tocdepth}{3}
\tableofcontents
}
\section{Overview}\label{overview}

The test programs in \texttt{vera-tests/} run on a fork of the
\textbf{atari800} emulator (version 5.2.0) that emulates the VeraX16
card on the Atari PBI bus: VERA registers at \texttt{\$D100-\$D11F}, the
2 KB handler ROM at \texttt{\$D800-\$DFFF} (selected through the PBI
latch \texttt{\$D1FF}).

The emulator lives next to this repository:

\begin{verbatim}
~/Progetti/ATARI/
├── atari800/                 emulator (src/atari800 is the binary)
│   └── vera_pbi_rom -> ../VERA_ATARI_PBI
└── VERA_ATARI_PBI/           this repository (ROM, drivers, tests)
\end{verbatim}

All the commands below are run \textbf{from the root of this repository}
(\texttt{VERA\_ATARI\_PBI}), so the emulator is
\texttt{../atari800/src/atari800}. For brevity:

\begin{Shaded}
\begin{Highlighting}[]
\VariableTok{EMU}\OperatorTok{=}\NormalTok{../atari800/src/atari800}
\end{Highlighting}
\end{Shaded}

\section{Prerequisites}\label{prerequisites}

\subsection{Building the emulator}\label{building-the-emulator}

\begin{Shaded}
\begin{Highlighting}[]
\BuiltInTok{cd}\NormalTok{ ../atari800}
\ExtensionTok{./autogen.sh}
\ExtensionTok{./configure} \AttributeTok{{-}{-}enable{-}pbi{-}verax16}
\FunctionTok{make}
\BuiltInTok{cd} \AttributeTok{{-}}
\end{Highlighting}
\end{Shaded}

The VERA emulation is in \texttt{src/pbi\_verax16.c} (registers, FX,
audio, SPI) and \texttt{src/vera\_video.c} (rendering). Without
\texttt{-\/-enable-pbi-verax16} the \texttt{-verax16} options do not
exist.

No Atari OS ROM is required: when none is configured in
\texttt{\textasciitilde{}/.atari800.cfg} atari800 uses the built-in
AltirraOS and Altirra BASIC.

\subsection{Building the ROM, the drivers and the test
disks}\label{building-the-rom-the-drivers-and-the-test-disks}

Required tools: \texttt{cc65} (\texttt{ca65}, \texttt{ld65},
\texttt{cl65}), \texttt{dir2atr}, \texttt{python3}.

\begin{Shaded}
\begin{Highlighting}[]
\FunctionTok{make}            \CommentTok{\# ROM, VERA*.SYS drivers, test programs, disk images}
\end{Highlighting}
\end{Shaded}

Main outputs:

\begin{longtable}[]{@{}
  >{\raggedright\arraybackslash}p{(\columnwidth - 2\tabcolsep) * \real{0.5000}}
  >{\raggedright\arraybackslash}p{(\columnwidth - 2\tabcolsep) * \real{0.5000}}@{}}
\toprule\noalign{}
\begin{minipage}[b]{\linewidth}\raggedright
File
\end{minipage} & \begin{minipage}[b]{\linewidth}\raggedright
Content
\end{minipage} \\
\midrule\noalign{}
\endhead
\bottomrule\noalign{}
\endlastfoot
\texttt{vera\_pbi\_handler.rom} & PBI handler ROM, 2 KB (boots in
80×60) \\
\texttt{VERA4030.SYS}, \texttt{VERA8030.SYS}, \texttt{VERA8060.SYS} &
RAM driver for 40×30, 80×30, 80×60 \\
\texttt{disk1-runcpm.atr} & DOS 2.0S, \texttt{RUNCPM.COM},
\texttt{VERA8030.SYS} \\
\texttt{disk2-veratests-40x30.atr} & \texttt{TEST4}, \texttt{TESTGS4},
\texttt{TESTMAZ4}, \texttt{TESTMTX4}, \texttt{TESTRMT},
\texttt{VERA4030.SYS} \\
\texttt{disk2-veratests-80x30.atr} & \texttt{TEST8}, \texttt{TESTGS8},
\texttt{TESTMAZ8}, \texttt{TESTMTX8}, \texttt{TESTRMT},
\texttt{VERA8030.SYS} \\
\texttt{disk2-veratests-80x60.atr} & \texttt{TEST6}, \texttt{TESTGS6},
\texttt{TESTMAZ6}, \texttt{TESTMTX6}, \texttt{TESTRMT},
\texttt{VERA8060.SYS} \\
\texttt{disk3-standalone.atr} & \texttt{TESTFX}, \texttt{TESTIRQ},
\texttt{TESTPLR} + \texttt{DEMO.VTM} (tests without the driver) \\
\texttt{disk4-rmtio.atr} & \texttt{TESTRIO} with MyPicoDos autorun and
the test files \\
\end{longtable}

The \texttt{TESTn}/\texttt{TESTGSn}/\texttt{TESTMAZn}/\texttt{TESTMTXn}
programs on the \texttt{disk2} disks already contain the driver of their
resolution (\texttt{n} = 4: 40×30, 8: 80×30, 6: 80×60).

\section{The basic command}\label{the-basic-command}

\begin{Shaded}
\begin{Highlighting}[]
\VariableTok{$EMU} \AttributeTok{{-}xl} \AttributeTok{{-}pal} \AttributeTok{{-}nopatch} \AttributeTok{{-}verax16} \AttributeTok{{-}verax16{-}rom}\NormalTok{ vera\_pbi\_handler.rom }\DataTypeTok{\textbackslash{}}
\NormalTok{     disk3{-}standalone.atr}
\end{Highlighting}
\end{Shaded}

\begin{itemize}
\tightlist
\item
  \texttt{-verax16} plugs the card in; \texttt{-verax16-rom} gives the
  handler ROM.
\item
  \texttt{-xl} (800XL) or \texttt{-xe} (130XE): the driver needs an
  XL/XE machine.
\item
  \texttt{-pal} or \texttt{-ntsc}: the tests must pass in both.
\item
  \texttt{-nopatch}: no high-speed SIO patch, disk I/O runs at the real
  serial speed (the H: device keeps working). To make it permanent, set
  \texttt{ENABLE\_SIO\_PATCH=0} in
  \texttt{\textasciitilde{}/.atari800.cfg}.
\item
  The last argument is the disk in \texttt{D1:}; more \texttt{.atr}
  files go to \texttt{D2:}, \texttt{D3:}\ldots{}
\end{itemize}

The emulator opens two outputs: the normal ANTIC/GTIA screen and the
VERA screen (640×480 VGA).

\subsection{Starting a program from DOS
2.0S}\label{starting-a-program-from-dos-2.0s}

The disks boot DOS 2.0S. From the DUP menu:

\begin{enumerate}
\def\labelenumi{\arabic{enumi}.}
\tightlist
\item
  press \texttt{L} (BINARY LOAD);
\item
  type the file name, for example \texttt{TESTFX.COM}, and press Return.
\end{enumerate}

\subsection{Starting a program without a
disk}\label{starting-a-program-without-a-disk}

\texttt{-run} loads an executable straight from the host file system:

\begin{Shaded}
\begin{Highlighting}[]
\VariableTok{$EMU} \AttributeTok{{-}xl} \AttributeTok{{-}pal} \AttributeTok{{-}nopatch} \AttributeTok{{-}verax16} \AttributeTok{{-}verax16{-}rom}\NormalTok{ vera\_pbi\_handler.rom }\DataTypeTok{\textbackslash{}}
     \AttributeTok{{-}run}\NormalTok{ TESTIRQ.COM}
\end{Highlighting}
\end{Shaded}

\section{Options}\label{options}

\subsection{VeraX16 card options}\label{verax16-card-options}

\begin{longtable}[]{@{}
  >{\raggedright\arraybackslash}p{(\columnwidth - 2\tabcolsep) * \real{0.5000}}
  >{\raggedright\arraybackslash}p{(\columnwidth - 2\tabcolsep) * \real{0.5000}}@{}}
\toprule\noalign{}
\begin{minipage}[b]{\linewidth}\raggedright
Option
\end{minipage} & \begin{minipage}[b]{\linewidth}\raggedright
Meaning
\end{minipage} \\
\midrule\noalign{}
\endhead
\bottomrule\noalign{}
\endlastfoot
\texttt{-verax16} & enable the card (alias \texttt{-\/-use-verax16}) \\
\texttt{-verax16-rom\ F} & handler ROM, 2 KB at
\texttt{\$D800-\$DFFF} \\
\texttt{-verax16-pbi-id\ N} & PBI device bit 0-7 (default 7, mask
\texttt{\$80}) \\
\texttt{-verax16-config-ms\ N} & FPGA configuration time after power-on,
bus dead meanwhile (default 100); \texttt{0} = the board holds Atari
RESET until CONFIG\_DONE \\
\texttt{-verax16-debuglevel\ N} & log level (default 0) \\
\texttt{-verax16-psg-volume\ N} & VERA PSG level in percent, 0-400
(default 100 = one voice as loud as a POKEY channel at volume 15) \\
\texttt{-verax16-sdcard\ F} & raw SD image (e.g.~from \texttt{dd})
exposed through the VERA SPI \\
\end{longtable}

\subsection{Useful atari800 options}\label{useful-atari800-options}

\begin{longtable}[]{@{}
  >{\raggedright\arraybackslash}p{(\columnwidth - 2\tabcolsep) * \real{0.5000}}
  >{\raggedright\arraybackslash}p{(\columnwidth - 2\tabcolsep) * \real{0.5000}}@{}}
\toprule\noalign{}
\begin{minipage}[b]{\linewidth}\raggedright
Option
\end{minipage} & \begin{minipage}[b]{\linewidth}\raggedright
Use
\end{minipage} \\
\midrule\noalign{}
\endhead
\bottomrule\noalign{}
\endlastfoot
\texttt{-xl}, \texttt{-xe} & 800XL / 130XE machine \\
\texttt{-pal}, \texttt{-ntsc} & video standard \\
\texttt{-basic}, \texttt{-nobasic} & BASIC ROM on/off \\
\texttt{-run\ F} & run a COM/EXE/XEX/BAS file \\
\texttt{-nopatch} & no high-speed SIO patch: real serial timing, H:
still works \\
\texttt{-nopatchall} & no OS patches: real SIO serial timing (H: does
not work) \\
\texttt{-volume\ N} & output volume 0-100 \\
\texttt{-stereo} & two POKEYs \\
\texttt{-turbo} & as fast as possible \\
\texttt{-netsio\ {[}port{]}} & NetSIO for FujiNet-PC (default UDP
9997) \\
\texttt{-config\ F} & alternate configuration file \\
\end{longtable}

\section{Running each test}\label{running-each-test}

\subsection{TESTFX --- FX coprocessor}\label{testfx-fx-coprocessor}

\begin{Shaded}
\begin{Highlighting}[]
\VariableTok{$EMU} \AttributeTok{{-}xl} \AttributeTok{{-}pal} \AttributeTok{{-}nopatch} \AttributeTok{{-}verax16} \AttributeTok{{-}verax16{-}rom}\NormalTok{ vera\_pbi\_handler.rom }\DataTypeTok{\textbackslash{}}
\NormalTok{     disk3{-}standalone.atr}
\end{Highlighting}
\end{Shaded}

From DUP: \texttt{L}, \texttt{TESTFX.COM}. It checks every FX register
and measures VRAM copy/fill throughput. Expected result: \textbf{PASS
36, FAIL 0}.

\subsection{TESTIRQ --- VERA interrupts}\label{testirq-vera-interrupts}

\begin{Shaded}
\begin{Highlighting}[]
\VariableTok{$EMU} \AttributeTok{{-}xl} \AttributeTok{{-}pal}  \AttributeTok{{-}nopatch} \AttributeTok{{-}verax16} \AttributeTok{{-}verax16{-}rom}\NormalTok{ vera\_pbi\_handler.rom }\DataTypeTok{\textbackslash{}}
     \AttributeTok{{-}run}\NormalTok{ TESTIRQ.COM}
\VariableTok{$EMU} \AttributeTok{{-}xl} \AttributeTok{{-}ntsc} \AttributeTok{{-}nopatch} \AttributeTok{{-}verax16} \AttributeTok{{-}verax16{-}rom}\NormalTok{ vera\_pbi\_handler.rom }\DataTypeTok{\textbackslash{}}
     \AttributeTok{{-}run}\NormalTok{ TESTIRQ.COM}
\end{Highlighting}
\end{Shaded}

It checks the \texttt{VIMIRQ} hook: VSYNC at 59.94 Hz, line IRQ, AFLOW
masking, chaining to the OS and removal. Expected result: \textbf{PASS
13, FAIL 0}, both in PAL and in NTSC. See
\texttt{Documentation/VERA-IRQ.md}.

\subsection{TESTPLR --- VTM player}\label{testplr-vtm-player}

\begin{Shaded}
\begin{Highlighting}[]
\VariableTok{$EMU} \AttributeTok{{-}xl} \AttributeTok{{-}pal} \AttributeTok{{-}nopatch} \AttributeTok{{-}volume}\NormalTok{ 100 }\AttributeTok{{-}verax16} \AttributeTok{{-}verax16{-}rom}\NormalTok{ vera\_pbi\_handler.rom }\DataTypeTok{\textbackslash{}}
\NormalTok{     disk3{-}standalone.atr}
\end{Highlighting}
\end{Shaded}

From DUP: \texttt{L}, \texttt{TESTPLR.COM}. Use \texttt{-pal} or
\texttt{-ntsc} according to the frame rate the song was converted for.
\texttt{-volume\ 100} is needed: the VERA PSG is quiet at the default
mixer level (or raise \texttt{-verax16-psg-volume}).

\subsection{TEST / TESTGS / TESTMAZ / TESTMTX --- driver and screen
modes}\label{test-testgs-testmaz-testmtx-driver-and-screen-modes}

\begin{Shaded}
\begin{Highlighting}[]
\VariableTok{$EMU} \AttributeTok{{-}xl} \AttributeTok{{-}pal} \AttributeTok{{-}nopatch} \AttributeTok{{-}verax16} \AttributeTok{{-}verax16{-}rom}\NormalTok{ vera\_pbi\_handler.rom }\DataTypeTok{\textbackslash{}}
\NormalTok{     disk2{-}veratests{-}80x30.atr}
\end{Highlighting}
\end{Shaded}

Use the disk of the resolution under test (\texttt{40x30},
\texttt{80x30}, \texttt{80x60}) and start, from DUP, \texttt{TEST8.COM},
\texttt{TESTGS8.COM}, \texttt{TESTMAZ8.COM}, \texttt{TESTMTX8.COM}
(digit 4 or 6 on the other disks): font loading, gradient, scrolling,
maze and matrix demos.

\subsection{TESTRMT --- RMT player on POKEY +
VERA}\label{testrmt-rmt-player-on-pokey-vera}

On every \texttt{disk2-veratests-*.atr} disk, it does not need
\texttt{VERA.SYS}. From DUP: \texttt{L}, \texttt{TESTRMT.COM}.

\begin{longtable}[]{@{}ll@{}}
\toprule\noalign{}
Key & Output \\
\midrule\noalign{}
\endhead
\bottomrule\noalign{}
\endlastfoot
\texttt{1} & POKEY only \\
\texttt{2} & VERA only \\
\texttt{3} & both (default) \\
\texttt{4} & hybrid, RMT8 only: POKEY channels 1-4, VERA channels 5-8 \\
\texttt{S} & VERA stereo on/off \\
\texttt{ESC} & stop and exit \\
\end{longtable}

\subsection{TESTRIO --- music during disk
I/O}\label{testrio-music-during-disk-io}

\begin{Shaded}
\begin{Highlighting}[]
\FunctionTok{make}\NormalTok{ disk4{-}rmtio.atr}
\VariableTok{$EMU} \AttributeTok{{-}xl} \AttributeTok{{-}pal} \AttributeTok{{-}nopatchall} \AttributeTok{{-}nobasic} \DataTypeTok{\textbackslash{}}
     \AttributeTok{{-}verax16} \AttributeTok{{-}verax16{-}rom}\NormalTok{ vera\_pbi\_handler.rom disk4{-}rmtio.atr}
\end{Highlighting}
\end{Shaded}

The disk starts \texttt{TESTRIO.COM} by itself (MyPicoDos autorun).
\texttt{-nopatchall} is mandatory: without it the emulator intercepts
SIO and there is no real serial timing. Expected (PAL, RBL loader):
\textasciitilde720 bytes/s, 0 errors, 0 lost ticks.

\subsection{RUNCPM --- ANSI terminal through
FujiNet}\label{runcpm-ansi-terminal-through-fujinet}

Terminal 1, FujiNet-PC:

\begin{Shaded}
\begin{Highlighting}[]
\BuiltInTok{cd}\NormalTok{ FujiNet/fujinet{-}pc{-}ATARI}
\ExtensionTok{./run{-}fujinet} \AttributeTok{{-}c}\NormalTok{ fnconfig.ini }\AttributeTok{{-}s}\NormalTok{ SD/}
\end{Highlighting}
\end{Shaded}

Wait for \texttt{\#\#\#\ NetSIO\ stopped\ \#\#\#} in the log, then in
terminal 2:

\begin{Shaded}
\begin{Highlighting}[]
\VariableTok{$EMU} \AttributeTok{{-}xl} \AttributeTok{{-}pal} \AttributeTok{{-}nopatch} \AttributeTok{{-}netsio} \AttributeTok{{-}verax16} \AttributeTok{{-}verax16{-}rom}\NormalTok{ vera\_pbi\_handler.rom}
\end{Highlighting}
\end{Shaded}

Details in \texttt{README-fujinet.md}.

\section{Troubleshooting}\label{troubleshooting}

\begin{longtable}[]{@{}
  >{\raggedright\arraybackslash}p{(\columnwidth - 2\tabcolsep) * \real{0.5000}}
  >{\raggedright\arraybackslash}p{(\columnwidth - 2\tabcolsep) * \real{0.5000}}@{}}
\toprule\noalign{}
\begin{minipage}[b]{\linewidth}\raggedright
Symptom
\end{minipage} & \begin{minipage}[b]{\linewidth}\raggedright
Cause / fix
\end{minipage} \\
\midrule\noalign{}
\endhead
\bottomrule\noalign{}
\endlastfoot
\texttt{ERROR:\ VeraX16\ PBI\ card\ not\ found\ (\$D100\ not\ responding).}
& \texttt{-verax16} is missing: every test calls
\texttt{vera\_require()} at start \\
Three beeps from the console speaker at boot & the ROM did not see the
VERA within \textasciitilde0.7 s; check \texttt{-verax16-config-ms} \\
Unknown option \texttt{-verax16} & emulator built without
\texttt{-\/-enable-pbi-verax16} \\
No sound or very quiet VERA & add \texttt{-volume\ 100} or raise
\texttt{-verax16-psg-volume} \\
A file is on the disk but DOS does not list it, or ``0 FREE SECTORS'' &
\texttt{dir2atr} bug; rebuild with \texttt{make}, which runs
\texttt{vera-tests/tools/fix\_atr\_vtoc.py} \\
SIO errors or wrong speed in \texttt{TESTRIO} & \texttt{-nopatchall} is
missing \\
\end{longtable}

\end{document}
